% Compléments pratiques après la bibliographie, sans modifier les quatre chapitres.
\clearpage
\phantomsection
\addcontentsline{toc}{chapter}{Annexes}
\chapter*{Mode d'emploi}
\addcontentsline{toc}{section}{Mode d'emploi}
\label{ann:demarrage}
Pour utiliser l'application, il faut ouvrir un terminal dans le dossier du projet, puis démarrer le serveur. Les commandes suivantes permettent de lancer la démonstration avec les modèles déjà entraînés.

\section*{Lancement de l'application}
Sous Windows, le script PowerShell permet de démarrer le serveur avec un environnement Python disponible sur l'ordinateur.
\begin{verbatim}
.\demo\start-reference.ps1
\end{verbatim}
On peut aussi installer les bibliothèques nécessaires et lancer le serveur directement avec Python.
\begin{verbatim}
python -m pip install -r requirements.txt
python -m demo.reference_server
\end{verbatim}
L'application s'ouvre dans le navigateur à l'adresse \url{http://127.0.0.1:8765/}. Le terminal doit rester ouvert pendant la démonstration. Pour arrêter le serveur, il suffit d'appuyer sur Ctrl+C dans ce terminal. Le lancement utilise les modèles enregistrés sans refaire l'entraînement.

\section*{Déroulement d'un essai}
Avant de lancer un essai, l'utilisateur choisit l'expérience, la forme, le nombre de robots et le contrôleur. Le numéro de départ permet de retrouver la même disposition initiale lorsque les autres réglages restent identiques.

La trajectoire est calculée avant le début de l'animation. Les robots passent d'abord par un cercle de rendez-vous, puis se répartissent les places et construisent la forme demandée. Le bouton Résultat final affiche directement le dernier état de la simulation.

Les commandes de lecture permettent de mettre l'animation en pause, de l'accélérer ou de revoir un instant précis. L'affichage des liens entre voisins, des trajectoires et des places aide à comprendre les mouvements. Les métriques présentées sous l'animation portent sur l'essai complet.

\chapter*{Les expériences}
\addcontentsline{toc}{section}{Les expériences}
\label{ann:experiences}
L'application propose sept expériences. Les choix disponibles dépendent du nombre de robots et du contrôleur utilisé. Les résultats présentés dans le mémoire concernent les configurations effectivement évaluées.

\section*{Formation et changement de forme}
\textbf{Former l'essaim.} Les robots construisent une lettre, un nombre ou une figure du catalogue après la préparation et l'attribution des places. L'erreur finale, les vitesses et les contacts permettent de juger le résultat. Le test final du GAT structuré utilise le rectangle, la lettre C et le nombre 3, avec 12, 24 et 64 robots.

\textbf{Changer de formation.} La deuxième étape reprend les positions et les vitesses obtenues à la fin de la première. Les places sont réattribuées pour la nouvelle forme. L'évaluation tient compte de chaque étape et des contacts sur toute la séquence. Cette expérience n'est pas proposée au modèle structuré ni à 64 robots dans la version étudiée.

\section*{Déplacement collectif}
\textbf{Déplacer en ligne droite.} Les positions cibles se déplacent ensemble le long d'un trajet sans obstacle. La forme doit être maintenue et le groupe doit atteindre la destination.

\textbf{Contourner un obstacle.} Le groupe suit un parcours qui passe à côté de l'obstacle. Le planificateur vérifie le passage de l'ensemble des places et adapte le détour si nécessaire.

\textbf{Obstacle sur la route.} Le trajet demandé est direct, mais un obstacle gêne certaines places. Le planificateur classique cherche un détour commun compatible avec la largeur de la formation.

\textbf{Naviguer avec plusieurs obstacles.} La scène comporte de cinq à dix obstacles proches des cibles. La campagne décrite dans le mémoire compare le contrôleur classique et le GAT initial protégé pour un carré de douze robots. Elle n'évalue pas le GAT structuré brut en navigation.

\textbf{Vérifier une destination obstruée.} Si une place finale est physiquement bloquée, la demande est signalée comme impossible. Ce signalement évite de lancer une mission dont la destination est inaccessible. Il reste distinct d'une formation réussie.

\chapter*{Résultats des simulations}
\addcontentsline{toc}{section}{Résultats des simulations}
\label{ann:resultats}
Les fichiers enregistrés permettent de retrouver les résultats présentés dans le mémoire et d'examiner les trajectoires. Les tests du programme et l'évaluation du modèle sont deux vérifications différentes.

\section*{Vérification du programme}
Les tests du logiciel s'exécutent depuis la racine du projet.
\begin{verbatim}
python -m pytest tests -q -p no:cacheprovider
\end{verbatim}
La commande suivante produit un rapport à partir des résultats déjà enregistrés. Un nouveau dossier de sortie permet de conserver les rapports précédents.
\begin{verbatim}
python -m demo.comparison_report --output outputs/lecture
\end{verbatim}

\section*{Fichiers des expériences}
Le rapport \texttt{008d-local-target-expanded-test.json} documente les neuf essais du GAT initial. Le rapport \texttt{014-vector-fresh-test.json} décrit les dix-huit essais réservés du modèle structuré. Le rapport \texttt{019-near-targets-evaluation.json} présente la campagne finale près des obstacles. Ces fichiers se trouvent dans \texttt{tasks/evidence/}. Les audits associés vérifient les trajectoires enregistrées.

Les journaux \texttt{run.json} des deux entraînements se trouvent dans leurs dossiers respectifs sous \texttt{outputs/imitation/}. Les poids conservés sont les fichiers \texttt{best.pt}. Le manifeste \texttt{latex/generated/sources.json} rassemble les empreintes des données utilisées pour les figures et les tableaux du mémoire.

\section*{Interprétation des résultats}
Pour comparer deux contrôleurs, on conserve la même forme, le même départ, le même effectif, les mêmes dimensions et les mêmes seuils. L'export contient les positions, les vitesses et les métriques. Une faible erreur finale ne suffit pas si un contact est survenu.

La stabilisation exige une erreur de formation au plus égale à 0,10~m et des vitesses au plus égales à 0,05~m/s pendant une seconde. La réussite exige aussi l'absence de contact sur l'essai. En navigation s'ajoutent l'arrivée à moins de 0,10~m et une déformation en route au plus égale à 0,20~m. Le temps indiqué est un temps simulé.
